\section{Maximum entropy problem}

%% Distributions
Given a non-negative and non-vanishing base measure $\basemeasurewith$, a probability distribution is a non-negative tensor $\probwith$ such that
\begin{align*}
    \contraction{\probwith,\basemeasurewith} = 1 \, .
\end{align*}
We denote the set of such distributions by
\begin{align*}
    \distmeasuredby{\basemeasure}
    = \big\{ \probwith \wcols &\probwith \geq \zerosat{\shortcatvariables} \ncond \contraction{\probwith,\basemeasurewith} = 1 \ncond \\
    &\forall_{\shortcatindicesin} \basemeasureat{\indexedshortcatvariables}=0 \Rightarrow \probat{\indexedshortcatvariables}=0  \big\} \, .
\end{align*}
Note that we here restrict the support of the tensor $\probwith$ to be included in the support of the base measure, to avoid ambiguity in the representation.

%% Entropy
The entropy of a distribution relative to the base measure $\basemeasure$ is
\begin{align*}
    \sentropyofwrt{\probwith}{\basemeasure}
    = -\contraction{\probwith,\lnof{\probwith},\basemeasurewith} \, .
\end{align*}
Note that this quantity depends on the chosen base measure.

%% Mean parameters
A statistic is a function $\sstat:\facstates\rightarrow\rr^\seldim$ assigning a real vector of dimension $\seldim$ to each state $\shortcatindicesin$.
We denote by $\sstatcoordinateof{\selindex}$ the restriction to the $\selindex$-th coordinate of $\sstat$.
The mean parameter of a distribution $\probwith$ to a statistic $\sstat$ is the vector $\meanparamwith\in\rr^\seldim$ with the coordinates
\begin{align*}
    \meanparamat{\indexedselvariable}
    = \expectationof{\enumformula}
    = \contraction{\sstatcoordinateofat{\selindex}{\shortcatvariables},\probwith,\basemeasurewith} \, .
\end{align*}
To express the computation of the mean parameter as a single contraction, we introduce a coordinate selection variable $\selvariable$ of dimension $\seldim$ and the selection encoding $\sencsstatwith$ with coordinates by
\begin{align*}
    \sencsstatat{\indexedshortcatvariables,\indexedselvariable}
    = \sstatcoordinateofat{\selindex}{\indexedshortcatvariables} \, .
\end{align*}
The mean parameter is then
\begin{align*}
    \meanparamwith
    = \contractionof{\probwith,\sencsstatwith}{\selvariable} \, .
\end{align*}

%% Maximum entropy problem
The maximum entropy problem given a mean parameter $\genmeanwith$ is stated by
\begin{equation}
    \tag{$\mathrm{P}_{\sstat,\meanparam,\basemeasure}$}\label{prob:maxEntropy}
    \argmax_{\probwith\in\distmeasuredby{\basemeasure}} \sentropyofwrt{\probwith}{\basemeasure}
    \stspace
    %\forall_{\selindexin} \wcols
    \contractionof{\probwith,\sencsstatwith,\basemeasurewith}{\selvariable} = \genmeanat{\selvariable}
\end{equation}